% Preámbulo común del material docente · XeLaTeX
\documentclass[11pt,a4paper]{article}
\usepackage[margin=22mm,top=20mm,bottom=22mm]{geometry}
\usepackage{fontspec}
\setmainfont{Avenir Next}[BoldFont={Avenir Next Demi Bold}, BoldItalicFont={Avenir Next Demi Bold Italic}]
\setsansfont{Avenir Next}
\setmonofont{Menlo}[Scale=0.88]
\newfontfamily\symfont{Apple Symbols}
\newfontfamily\unifont{Arial Unicode MS}
\usepackage{unicode-math}
\setmathfont{latinmodern-math.otf}
\usepackage{polyglossia}\setdefaultlanguage{spanish}
\usepackage{xcolor,booktabs,tabularx,enumitem,parskip,microtype,titlesec,fancyhdr}
\usepackage[most]{tcolorbox}
\usepackage[hidelinks]{hyperref}
\emergencystretch=2em

\definecolor{tinta}{HTML}{1F2933}
\definecolor{acento}{HTML}{B7410E}
\definecolor{suave}{HTML}{F4F1EC}
\definecolor{gris}{HTML}{6B7280}
\definecolor{linea}{HTML}{D6D0C6}
\color{tinta}

\titleformat{\section}{\Large\bfseries\color{acento}}{\thesection}{0.6em}{}
\titleformat{\subsection}{\large\bfseries}{\thesubsection}{0.6em}{}
\titlespacing*{\section}{0pt}{1.6em}{0.5em}
\titlespacing*{\subsection}{0pt}{1.1em}{0.3em}
\setlist{nosep, leftmargin=1.4em, itemsep=2pt}
\renewcommand{\arraystretch}{1.25}

% bloque de órdenes de terminal
\newtcblisting{orden}{listing only, colback=suave, colframe=linea, boxrule=0.4pt, arc=2pt,
  left=6pt, right=6pt, top=4pt, bottom=4pt, listing options={basicstyle=\ttfamily\small, breaklines=true, columns=fullflexible, keepspaces=true}}
% nota destacada
\newtcolorbox{nota}[1][Nota]{colback=white, colframe=acento, boxrule=0.6pt, arc=2pt, title=#1,
  fonttitle=\bfseries\small, coltitle=white, colbacktitle=acento, left=6pt, right=6pt}
% preguntas
\newcommand{\pregunta}[2]{\item[\textbf{(#1)}] #2}
\newenvironment{preguntas}{\begin{description}[leftmargin=2.4em, style=nextline, itemsep=4pt]}{\end{description}}
\newcommand{\ok}{{\symfont ✔}}
\newcommand{\nok}{{\symfont ✘}}
\newcommand{\qed}{{\symfont ∎}}
\newcommand{\codigo}[1]{\texttt{#1}}

\pagestyle{fancy}\fancyhf{}
\renewcommand{\headrulewidth}{0pt}
\fancyfoot[C]{\small\color{gris}\docpie\ · página \thepage}

\newcommand{\cabecera}[3]{%
  {\Huge\bfseries\color{acento}#1\par}\vspace{4pt}
  {\large #2\par}\vspace{2pt}
  {\color{gris}#3\par}\vspace{8pt}
  {\color{linea}\hrule height 0.6pt}\vspace{10pt}}

\newcommand{\docpie}{Laboratorio 1 · teoría mínima de RL}
\usepackage{listings}
\lstset{basicstyle=\ttfamily\small, language=Python, columns=fullflexible, keepspaces=true, breaklines=true,
  keywordstyle=\color{acento}, commentstyle=\color{gris}\itshape, stringstyle=\color{tinta},
  showstringspaces=false, frame=none, xleftmargin=0pt, inputencoding=utf8, extendedchars=true}
\newtcblisting{python}{listing only, colback=suave, colframe=linea, boxrule=0.4pt, arc=2pt,
  left=6pt, right=6pt, top=4pt, bottom=4pt, listing options={style=tcblatex, basicstyle=\ttfamily\small,
  language=Python, columns=fullflexible, keepspaces=true, breaklines=true, keywordstyle=\color{acento},
  commentstyle=\color{gris}\itshape, showstringspaces=false}}
\begin{document}
\cabecera{Teoría mínima de RL}{Lo que hace falta saber para el laboratorio 1, y cómo está escrito en el código}{Métodos numéricos · Matemáticas · lectura de 20 minutos · no hay que memorizar nada}

\section{La idea en una frase}

\textbf{Aprendizaje por refuerzo} (\emph{reinforcement learning}, RL): un programa toma decisiones, recibe puntos por sus consecuencias, y ajusta sus decisiones para ganar más puntos. Nadie le dice qué hacer; solo cuánto valió lo que hizo.

Es distinto de programar el método: en vez de escribir «corta en el medio», escribimos «te doy puntos por encoger el intervalo y te los quito por perder la raíz», y miramos si «cortar en el medio» \emph{emerge}.

\section{Los cinco ingredientes}

\begin{tabularx}{\linewidth}{@{}l X X@{}}
\toprule
\textbf{nombre} & \textbf{qué es} & \textbf{en la bisección} \\
\midrule
\textbf{agente} & quien decide & el programa que no conoce el método \\
\textbf{entorno} & el mundo donde decide; responde a cada decisión con una observación y puntos & una función $f$ continua y un intervalo $[a,b]$ con $f(a)f(b)<0$ \\
\textbf{estado} $s$ & lo que el agente \emph{ve} antes de decidir & los signos de $f(a)$, $f(x)$, $f(b)$ \\
\textbf{acción} $a$ & lo que puede hacer & la fracción $\lambda$ donde cortar; qué lado conservar \\
\textbf{recompensa} $r$ & los puntos que recibe justo después de actuar & $-1 + \log_2\dfrac{\text{ancho anterior}}{\text{ancho nuevo}}$; $-10$ si pierde la raíz \\
\bottomrule
\end{tabularx}

Un \textbf{episodio} (en la interfaz, \emph{generación}) es una partida completa: desde el intervalo inicial hasta que el ancho baja de la tolerancia, o hasta que se pierde la raíz. Se juegan miles de partidas, cada una con una $f$ distinta.

Una \textbf{política} $\pi$ es la regla del agente: en el estado $s$, hacer la acción $\pi(s)$. Al principio es aleatoria. El objetivo es encontrar la política que maximiza el \textbf{retorno}: la suma de recompensas del episodio,
\[
G = r_1 + r_2 + \dots + r_n .
\]
(En general se usa $G = \sum_k \gamma^{k} r_k$ con un factor de descuento $\gamma\in[0,1]$ que da menos peso al futuro. En este laboratorio no hace falta: la bisección usa $\gamma = 0$, y la demostración $\gamma = 0{,}97$.)

\begin{nota}[Por qué el retorno de la bisección mide lo correcto]
La recompensa de progreso \emph{telescopa}: si $w_0$ es el ancho inicial y $w_n$ el final,
\[
\sum_{k=1}^{n}\Bigl(-1+\log_2\tfrac{w_{k-1}}{w_k}\Bigr) = \log_2\tfrac{w_0}{w_n} - n .
\]
El primer sumando está fijado por la tolerancia. Así que maximizar el retorno es exactamente \textbf{minimizar $n$}, el número de evaluaciones de $f$. La recompensa paso a paso solo sirve para que el agente vea progreso antes de terminar; no cambia qué es óptimo. (Esto se llama \emph{shaping basado en potencial}, Ng, Harada y Russell, 1999.)
\end{nota}

\section{La tabla Q}

El agente guarda una tabla con una fila por estado y una columna por acción. La entrada $Q(s,a)$ es su estimación de \textbf{cuántos puntos va a ganar en total} si en el estado $s$ hace la acción $a$ y después sigue actuando bien. Al principio toda la tabla vale 0.

\begin{center}
\begin{tabular}{l|ccccc}
 & $\lambda=0{,}1$ & $\lambda=0{,}2$ & $\dots$ & $\lambda=0{,}5$ & $\dots$ \\
\hline
estado «corte» & $-0{,}55$ & $-0{,}28$ & & $\mathbf{0{,}00}$ & \\
\end{tabular}
\qquad
\begin{tabular}{l|cc}
 & izq & der \\
\hline
$(+,-,+)$ & $\mathbf{0{,}0}$ & $-10$ \\
$(+,+,-)$ & $-10$ & $\mathbf{0{,}0}$ \\
$(-,+,+)$ & $\mathbf{0{,}0}$ & $-10$ \\
$(-,-,+)$ & $-10$ & $\mathbf{0{,}0}$ \\
\end{tabular}
\end{center}

Con la tabla aprendida, la política es leerla: en cada estado, la acción con mayor $Q$. Eso es la \textbf{política voraz} (\emph{greedy}). En el informe, la tabla de la izquierda es «política de corte» y la de la derecha «regla de conservación».

\subsection*{Cómo se actualiza: Q-learning}

Después de cada acción el agente tiene cuatro datos: estaba en $s$, hizo $a$, recibió $r$ y llegó a $s'$. Entonces corrige la entrada:
\[
Q(s,a) \;\leftarrow\; Q(s,a) + \alpha\,\Bigl[\,\underbrace{r + \gamma \max_{a'} Q(s',a')}_{\text{lo que ahora cree que vale}} - Q(s,a)\Bigr].
\]
Es decir: mueve la estimación vieja una fracción $\alpha$ hacia la nueva evidencia. La nueva evidencia es «lo que gané ahora» más «lo mejor que creo poder ganar desde donde caí». Con $\gamma=0$ el segundo término desaparece y $Q(s,a)$ es simplemente el \textbf{promedio de las recompensas} obtenidas al hacer $a$ en $s$; con $\alpha = 1/N$ (siendo $N$ las veces que se ha probado) ese promedio es exacto.

\textbf{Ejemplo a mano.} Estado «corte», acción $\lambda = 0{,}3$, $Q = 0$ al principio. Primera vez: la raíz cae en el 30\,\% izquierdo, se conserva ese lado, progreso $\log_2(1/0{,}3) = 1{,}74$, recompensa $r = 0{,}74$. Con $\alpha = 1/1$: $Q \leftarrow 0{,}74$. Segunda vez: la raíz cae a la derecha, progreso $\log_2(1/0{,}7)=0{,}51$, $r=-0{,}49$. Con $\alpha=1/2$: $Q \leftarrow 0{,}74 + \tfrac12(-0{,}49-0{,}74) = 0{,}125$. Tras cientos de veces, $Q(0{,}3)$ converge a su promedio, unos $-0{,}12$. Y $Q(0{,}5)$ converge a $0$, porque con $\lambda = 1/2$ el progreso es siempre exactamente 1 bit. Por eso gana.

\subsection*{Explorar o explotar}

Si el agente siempre hiciera lo que la tabla dice mejor, nunca probaría las otras acciones y no descubriría que son mejores. Se resuelve con \textbf{$\varepsilon$-greedy}: con probabilidad $\varepsilon$ hace una acción al azar (explora); si no, la mejor según la tabla (explota). $\varepsilon$ empieza en 1 (todo al azar) y baja poco a poco; en el marcador de la interfaz se ve como «ε exploración». Además hay un \textbf{bono de curiosidad} (UCB): a la hora de elegir, cada acción suma $c\sqrt{\ln N / (n_a+1)}$, mayor cuanto menos se ha probado, para que las acciones poco visitadas parezcan atractivas hasta que se conozcan bien.

Esto explica dos cosas del laboratorio: por qué los hitos ocurren en generaciones distintas con otra semilla (la exploración es aleatoria), y por qué «la política voraz ya produce la demostración mínima» puede ocurrir \emph{antes} de que el agente la produzca entrenando (mientras entrena, sigue explorando).

\section{Dos agentes en la fase 1, y por qué}

En la bisección hay dos decisiones por paso y se usan dos tablas independientes (dos «cabezas»):
\begin{itemize}
\item \textbf{cabeza de corte:} un solo estado y 9 acciones ($\lambda = 0{,}1,\dots,0{,}9$). La fracción óptima no depende del ancho ni de la $f$, así que todas las muestras refuerzan la misma fila. Es un \emph{bandido de 9 brazos}.
\item \textbf{cabeza de conservación:} estado = signos de $(f(a), f(x), f(b))$, 2 acciones. Solo los signos deciden dónde está la raíz, y son lo único que el método real puede ver.
\end{itemize}
Un detalle de diseño: si se pierde la raíz, la culpa es de la conservación, no del corte, así que ese $-10$ solo actualiza la segunda tabla.

\section{La fase 2: la demostración como juego}

Mismo esquema, otro entorno. El \textbf{estado} es el conjunto de pasos ya establecidos (una lista de 21 casillas: 14 pasos verdaderos y 7 distractores). La \textbf{acción} es afirmar el siguiente paso. El entorno tiene un \textbf{verificador}: un paso es válido si todas sus dependencias ya están establecidas, no se ha usado, y no es distractor. \textbf{Recompensa}: $-0{,}5$ por paso válido (cada línea cuesta), $-2$ por paso inválido («salto lógico»), $+20$ al llegar a $\qed$. Aquí sí importa el futuro ($\gamma = 0{,}97$): el valor del $+20$ se propaga hacia atrás por la cadena de dependencias, y por eso los pasos profundos entran en la política al final.

\clearpage
\section{Cómo está escrito en el código}

Todo lo anterior son unas 150 líneas en \codigo{01\_biseccion/bisectrl/}. Estas son las que definen el juego.

\subsection*{Las funciones de prueba (\codigo{env\_bisection.py}, \codigo{make\_problem})}

En cada episodio se sortea una familia y unos parámetros. Todas tienen cambio de signo en $[a,b]$ y el entorno calcula la raíz de referencia con una bisección exacta (que el agente nunca ve).

\begin{tabularx}{\linewidth}{@{}l X X@{}}
\toprule
\textbf{familia} & \textbf{fórmula} & \textbf{parámetros aleatorios} \\
\midrule
\codigo{poly3} & $(x-r)\bigl((x-p)^2+q\bigr)$ & $r\in[-1{,}5, 1{,}5]$, $p\in[-2,2]$, $q\in[0{,}2,2]$; una sola raíz real \\
\codigo{sin} & $\sin(kx)-c$ & $k\in[0{,}7,1{,}6]$, $c\in[-0{,}8,0{,}8]$ \\
\codigo{exp} & $e^{x}-c$ & $c\in[1{,}5,6]$ \\
\codigo{cos\_mix} & $\cos x - x$ & solo el intervalo: $a\in[-1,0{,}3]$, $b\in[1{,}2,2{,}5]$ \\
\codigo{cubic\_root} & $(x-r)^3+0{,}3(x-r)$ & $r\in[-1,1]$; casi plana cerca de la raíz \\
\bottomrule
\end{tabularx}

\begin{python}
def make_problem(rng):
    kind = rng.choice(["poly3", "sin", "exp", "cos_mix", "cubic_root"])
    if kind == "poly3":
        r = rng.uniform(-1.5, 1.5)
        p, q = rng.uniform(-2, 2), rng.uniform(0.2, 2)
        f = lambda x, r=r, p=p, q=q: (x - r) * ((x - p) ** 2 + q)
        a, b = r - rng.uniform(0.5, 2.5), r + rng.uniform(0.5, 2.5)
    ...
    if f(a) * f(b) >= 0:          # por seguridad numérica
        return make_problem(rng)
    root = _bisect_reference(f, a, b)
    return Problem(name, f, a, b, root)
\end{python}

Para añadir una familia basta con otro \codigo{elif} que defina \codigo{f}, \codigo{a} y \codigo{b}. Es un buen primer cambio que hacer.

\subsection*{El estado: qué ve el agente}
\begin{python}
def state_split(self):
    # un unico estado: la fraccion optima no depende del ancho
    return ("split",)

def state_keep(self):
    # solo los signos: es lo unico que decide que lado conserva la raiz
    sa = 1 if self.fa > 0 else -1
    sx = 1 if self.fx > 0 else (-1 if self.fx < 0 else 0)
    sb = 1 if self.fb > 0 else -1
    return ("keep", sa, sx, sb)
\end{python}
Fíjate en lo que \textbf{no} está en el estado: ni el valor de $f$, ni el ancho, ni la raíz. Si el estado contuviera la raíz, el agente «aprendería» a mirarla. La regla es: solo cosas que el método real puede medir.

\subsection*{La recompensa: qué le da puntos}
\begin{python}
REWARD_STEP = -1.0        # cada evaluacion de f cuesta
REWARD_LOST_ROOT = -10.0  # conservar el lado sin cambio de signo

def step_keep(self, keep):
    old_width = self.b - self.a
    ...  # actualiza [a, b] al lado elegido
    if self.fa * self.fb > 0:
        return self.state_split(), REWARD_LOST_ROOT, True, "lost"
    progress = math.log2(old_width / (self.b - self.a))
    if self.b - self.a < self.tol:
        return self.state_split(), REWARD_STEP + progress, True, "converged"
    return self.state_split(), REWARD_STEP + progress, False, "step"
\end{python}
Cada \codigo{return} devuelve cuatro cosas: el nuevo estado, la recompensa, si el episodio terminó, y una etiqueta para la interfaz. En el laboratorio cambiarás \codigo{REWARD\_LOST\_ROOT} a 0 y verás qué pasa. La recompensa es la especificación completa del problema: si está mal, el agente aprende otra cosa, y lo aprende bien.

\subsection*{El agente: la tabla y su actualización (\codigo{agent.py})}
\begin{python}
def act(self, state, actions=None):
    if self.rng.random() < self.eps:          # explorar
        return self.rng.choice(actions)
    if self.ucb_c > 0:                        # bono de curiosidad
        qs, ns = self.q[state], self.counts[state]
        n_state = sum(ns) + 1
        score = {a: qs[a] + self.ucb_c * math.sqrt(math.log(n_state) / (ns[a] + 1)) for a in actions}
        return max(actions, key=score.get)
    return self.greedy(state, actions)        # explotar

def learn(self, s, a, r, s2, done):
    target = r if done else r + self.gamma * max(self.q[s2])
    self.counts[s][a] += 1
    alpha = self.alpha if self.alpha is not None else 1.0 / self.counts[s][a]
    self.q[s][a] += alpha * (target - self.q[s][a])
\end{python}
La línea de \codigo{learn} que empieza por \codigo{self.q[s][a] +=} es exactamente la fórmula de la sección 3. Es todo el «aprendizaje».

\subsection*{El verificador de la demostración (\codigo{proof\_kb.py}, \codigo{env\_proof.py})}
\begin{python}
Step("ROOT", "Raíz",
     "El límite de una sucesión <= 0 es <= 0: f(c)^2 <= 0, así que f(c) = 0.",
     deps=("INV", "CONT"))

def is_valid(self, action, mask=None):
    step = STEPS[action]
    if step.distractor or (mask >> action) & 1:   # distractor o repetido
        return False
    return (DEPS_MASK[action] & mask) == DEPS_MASK[action]   # todas las deps establecidas
\end{python}
El estado es un entero cuyos bits dicen qué pasos están establecidos; \codigo{is\_valid} comprueba con una operación de bits que todas las dependencias están. Eso es todo lo que el «verificador» sabe de matemáticas: \textbf{lo que dicen las \codigo{deps}}. Si en \codigo{ROOT} quitas \codigo{"INV"} de \codigo{deps}, el verificador acepta con gusto una demostración que no prueba el invariante. Ese es el experimento \codigo{grafo} del laboratorio.

\section{Vocabulario para la sesión}

\begin{tabularx}{\linewidth}{@{}l X@{}}
\toprule
episodio / generación & una partida completa, con una $f$ nueva \\
retorno & suma de recompensas del episodio \\
política & regla estado $\to$ acción; \emph{voraz} = leer el máximo de la tabla $Q$ \\
$Q(s,a)$ & estimación del retorno futuro tras hacer $a$ en $s$ \\
$\alpha$ & tasa de aprendizaje: cuánto se mueve $Q$ hacia la nueva evidencia \\
$\gamma$ & descuento: cuánto pesa el futuro (0: nada; 1: igual que el presente) \\
$\varepsilon$ & probabilidad de explorar al azar \\
shaping & recompensa intermedia que mide progreso sin cambiar qué es óptimo \\
distractor & paso de la demostración que nunca es válido: falso, irrelevante o non sequitur \\
salto lógico & afirmar un paso cuyas dependencias no están establecidas \\
\bottomrule
\end{tabularx}

\vspace{6pt}
\textbf{Para leer más:} Sutton y Barto, \emph{Reinforcement Learning: An Introduction} (2.ª ed., 2018), capítulos 2 (bandidos) y 6 (Q-learning). Gratis en línea. Los capítulos son más fáciles de lo que parecen.
\end{document}
